\chapterrule
\section{怎样改进回答：训练反馈与推理时计算}
\label{intro:23c-post-training}
\label{intro:23cc-test-time-compute}

\begin{idea}
\noindent\emph{任务目标改变，训练信号也需要改变}\par

报告助手要学会完整输出单位与出处，并在缺证据时指出缺口。
\textbf{任务完成质量}要求清楚的目标和反馈：
遵守格式、核对证据、处理不确定性。
这里比较监督微调、偏好学习、可验证奖励和推理时策略
如何分别改变模型或生成过程。
\end{idea}

\subsection{先看训练与推理策略怎样分工}
\phantomsection\label{sec:training-inference-roles}

先只问三个问题：反馈来自哪里，是否更新模型参数，额外计算发生在训练还是当前请求？
后面的各节按这张表展开：

\begin{center}
\small
\renewcommand{\arraystretch}{1.25}
\tagpdfsetup{table/header-rows={1}}
\begin{tabularx}{\linewidth}{@{}p{29mm}p{36mm}p{19mm}X@{}}\toprule
方法 & 反馈或输入信号 & 更新参数？ & 主要改变什么\\\midrule
预训练 & 文本中的后续 Token & 是 & 学习一般的语言模式\\
示范答案训练 & 人工示范答案 & 是 & 学习特定任务的回答方式\\
偏好训练 & 回答之间的偏好比较 & 是 & 让较好的回答更容易生成\\
程序核验训练 & 程序检查结果给出的奖励 & 是 & 按可验证结果改进回答策略\\
步骤提示 & 当前输入中的解题步骤示例 & 否 & 改变本次输入的组织方式\\
多份候选再筛选 & 当前请求中的生成与核验 & 否 & 把更多计算用于筛选本次回答\\
\bottomrule
\end{tabularx}
\end{center}

因此，本章会先讲训练怎样改变参数，再讲固定参数时怎样把更多预算用在一次回答上。

\subsection{预训练与监督微调：看起来都是预测 Token，但哪一项变了？}
\phantomsection\label{concept:sft}

预训练使用文本中的后续 Token 作为目标；
监督微调（SFT）使用特定任务的输入和示范答案。
数学上，两者都可以使用 Token 级 \conceptlink{cross-entropy}{交叉熵}，
但数据如何构造、哪些位置计入损失不同。

例如训练对话数据为：
\[
\text{用户：$2+3=?$}\qquad
\text{助手：$5$}.
\]
常见做法是把用户提问作为条件，
重点让助手回答的 Token 参与损失；
也可以按实际训练目标选择不同掩码。
这样模型不仅学语言续写，还学习
在给定指令和角色约束下如何回应。

\subsubsection*{聊天模板是角色进入 Token 序列的接口}
\phantomsection\label{concept:chat-template}
应用保存的 \texttt{role=user/assistant/system} 消息列表，
要先由模型配套的\textbf{聊天模板}序列化，再送进 Tokenizer／模型。
例如一种示意格式是
\begin{quote}\ttfamily\small
<user> 2+3=? <end> <assistant> 5 <end>
\end{quote}
真实模板的角色标记、分隔符、回合结束符与专用 ID 由该模型定义；
有的格式由多个普通 Token 拼成，有的标记对应单个特殊 Token。
生成前补上助手前缀，模型才从正确角色的位置续写。
训练和推理应使用匹配模板，避免重复添加 BOS/EOS；
SFT 的损失掩码再决定哪些回答位置计分。
这些标记说明消息结构，权限仍由应用和工具执行层控制。

\tagstructbegin{tag=Figure,alttext={同一份季度报告格式示范的两条路径：放入本次 Prompt 时，只参与前向计算，模型参数保持不变；作为 SFT 训练样本时，计算交叉熵与梯度，优化器更新参数。原始报告事实还要由文件读取和校验提供。}}
\begin{center}
\begin{tikzpicture}[>=Latex,node distance=6mm,
 n/.style={draw=linegray,rounded corners=2pt,fill=paper,align=center,
 font=\small,text width=33mm,minimum height=14mm}]
\node[n] (data) {报告格式示范\\季度、收入、出处};
\node[n,above right=3mm and 9mm of data,fill=blueback] (prompt) {放入 Prompt\\本次上下文};
\node[n,right=of prompt] (fixed) {模型前向\\$\theta$ 保持不变};
\node[n,below right=3mm and 9mm of data,fill=orangeback] (sft) {构造 SFT 样本\\答案 Token 计损失};
\node[n,right=of sft] (update) {反传＋优化器\\$\theta\to\theta'$};
\draw[->,thick,blue] (data)--(prompt);
\draw[->,thick,blue] (prompt)--(fixed);
\draw[->,thick,orange] (data)--(sft);
\draw[->,thick,orange] (sft)--(update);
\end{tikzpicture}
\captionof{figure}{同一示范放在哪里，决定它改变当前上下文还是训练参数}
\label{fig:training-context-branches}
\end{center}
\tagstructend


\subsection{偏好学习：一个问题怎样选择更有用的回答}
\phantomsection\label{concept:rlhf}

面对“请解释什么是梯度”，设回答 A 准确、循序渐进，
回答 B 虽流畅却有明显错误。
两种回答都可以由语言模型生成，
我们希望训练后，模型更容易生成 A 这样的回答。
先把候选简化成 A、B 两项：若原来各有 $0.5$ 的概率，更新后变成 $(0.7,0.3)$，
较好回答就更容易被选中；这些新概率要由训练得到。

给定问题后，模型对各个回答分配的生成概率，称为\textbf{策略}，记作 $\pi_\theta(y\mid x)$。
这里 $x$ 是问题、$y$ 是回答，$\theta$ 是模型参数；整条回答的概率来自各 Token 条件概率的乘积。
因此\textbf{策略优化就是根据反馈更新参数，改变以后生成回答的概率}。
基于人类反馈的强化学习（RLHF）利用人的偏好或评价提供这种反馈。
经典流程先训练\phantomsection\label{concept:reward-model}\textbf{奖励模型} $r_\phi(x,y)$，
让它依据偏好数据给回答评分，再用这些分数优化语言模型。
下面先算一条成对偏好怎样训练奖励模型。

\subsection{成对偏好怎样训练奖励模型？}
把问题 $x$ 下偏好的回答 A 与较差回答 B 的分数设为
$r_\phi(x,A)=1$、$r_\phi(x,B)=0$。
先用 sigmoid 把任意实数转成 $0$ 到 $1$ 的值：
$\operatorname{sigmoid}(u)=1/(1+e^{-u})$。
Bradley--Terry 型成对模型把分数差变成偏好概率：
\[
P(A\succ B\mid x)=\operatorname{sigmoid}\bigl(r_\phi(x,A)-r_\phi(x,B)\bigr)
=\frac{1}{1+e^{-1}}\approx\boxed{0.731}.
\]
若偏好标签是“A 优于 B”，负对数损失为
\[
-\ln P(A\succ B)\approx\boxed{0.313}.
\]
这项损失让奖励模型参数 $\phi$ 学会把偏好的回答评分更高；实际训练会用许多成对标注共同确定 $\phi$。
偏好数据质量要通过标注一致性、分组评测及失败案例验证。

奖励模型提供优化方向，但单追高分也可能让当前策略偏离原来有用的语言行为。
常见做法是保留一份冻结的 SFT 模型作为\textbf{参考策略} $\pi_{\rm ref}$，
用 KL 散度衡量当前策略与参考策略的回答分布差异，并把它计入目标：
\[
\max_\theta\quad
\mathbb E_{x\sim\mathcal D,\,y\sim\pi_\theta(\cdot\mid x)}[r_\phi(x,y)]
-\beta_{\rm KL}\,\mathbb E_{x\sim\mathcal D}
\left[D_{\rm KL}\bigl(\pi_\theta(\cdot\mid x)\Vert\pi_{\rm ref}(\cdot\mid x)\bigr)\right].
\]
KL 惩罚让提高奖励的同时也付出偏离参考分布的代价。仍以参考概率 $(0.5,0.5)$、新概率 $(0.7,0.3)$ 为例：
\[
D_{\rm KL}(\pi_\theta\Vert\pi_{\rm ref})
=0.7\ln\frac{0.7}{0.5}+0.3\ln\frac{0.3}{0.5}\approx0.0823.
\]
两份分布相同时此值为零；本例的 $0.0823$ 量出了这次概率调整的偏离程度。
下文的直接偏好优化（DPO）把偏好损失直接写到模型概率上，跳过单独的奖励模型。

\subsubsection*{\texorpdfstring{{\symboliconfont ★}}{★} 选读 DPO：让对数概率比充当隐式奖励}
\phantomsection\label{concept:dpo}
\textbf{直接偏好优化（Direct Preference Optimization, DPO）}把同一个成对损失直接写到语言模型上。
对偏好回答 $y_w$ 与较差回答 $y_l$，各算一次当前策略相对参考策略的对数概率比
\[
\Delta_w=\ln\frac{\pi_\theta(y_w\mid x)}{\pi_{\rm ref}(y_w\mid x)},\qquad
\Delta_l=\ln\frac{\pi_\theta(y_l\mid x)}{\pi_{\rm ref}(y_l\mid x)},\qquad
\ell_{\rm DPO}=-\ln\operatorname{sigmoid}\bigl(\beta_{\rm DPO}(\Delta_w-\Delta_l)\bigr).
\]
为避免与采样章的逆温度 $\beta=1/\tau$ 混淆，本节把 DPO 的系数记为 $\beta_{\rm DPO}$。
$\beta_{\rm DPO}\Delta$ 扮演\textbf{隐式奖励}：在带 KL 约束的 RLHF 目标下，最优策略满足
$r(x,y)=\beta_{\rm DPO}\ln\bigl(\pi(y\mid x)/\pi_{\rm ref}(y\mid x)\bigr)+\text{只依赖 }x\text{ 的常数}$，
成对相减时常数抵消。所以 DPO 不必另训奖励模型，训练时也不必从模型采样。
这里的两个回答处在完整回答空间中，其他回答也占有概率质量；不能把这两个比值当成仅含 $y_w,y_l$ 的二项分布。
取 $\beta_{\rm DPO}=1$，若当前策略相对参考策略把 $y_w$ 的对数概率提高 $0.5$、把 $y_l$ 的降低 $0.5$，
则 $\Delta_w-\Delta_l=1$，
\[
\ell_{\rm DPO}=-\ln\operatorname{sigmoid}(1)\approx\boxed{0.313},
\]
与上面 $r_A-r_B=1$ 的成对损失相同。继续训练会增大这个差，$\beta_{\rm DPO}$ 越大，同样的概率比变化被放大得越多。
公式见 \href{https://arxiv.org/abs/2305.18290}{Rafailov 等（2023）的 DPO 论文}。

\subsection{可验证奖励：让模型在能自动检查的任务中尝试改进}
\phantomsection\label{concept:rlvr}

对于一部分题目，我们可以编写\textbf{自动核验结果}的程序。
例如问题是计算 $17+25$，提交答案 $42$ 可以得奖励 $1$，
提交其他答案得 $0$。
这类信号可以用于基于可验证奖励的强化学习（常见缩写 RLVR）。
RLVR 说明奖励来自可核验结果；下一步还要选算法，把奖励变成参数更新。

\begin{center}
\begin{tikzpicture}[>=Latex,node distance=5mm,
 n/.style={draw=linegray,rounded corners=3pt,fill=paper,minimum width=108mm,minimum height=10mm,align=center,font=\small}]
\node[n] (q) {问题：$17+25$；模型采样一个候选答案};
\node[n,below=of q] (v) {验证器执行规则：答案是否为 $42$？};
\node[n,below=of v] (r) {奖励信号 $\to$ 学习算法更新模型策略};
\draw[->,thick,flowcolor] (q)--(v);
\draw[->,thick,flowcolor] (v)--(r);
\end{tikzpicture}
\end{center}

\begin{example}
\noindent\emph{算一组可验证奖励与相对优势}\par
若模型对同一道题独立生成四个答案 $(42,41,42,50)$，验证器给予奖励 $(1,0,1,0)$。
怎样把奖励变成“这次回答应该更常出现，还是更少出现”的信号？
先与同组的平均表现比较。本组平均奖励为 $\bar r=0.5$，
答对的偏差为 $1-0.5=+0.5$，答错的为 $0-0.5=-0.5$。
这样正号表示比同组平均好，负号表示比同组平均差。

\phantomsection\label{concept:grpo}
\textbf{组相对策略优化（GRPO）}使用同一问题的一组候选奖励构造相对优势，
再据此优化策略。这里先计算它所需的优势信号。
GRPO 指更新算法，奖励既可以来自程序验证器，也可以来自奖励模型；
本例选择程序核验 $42$ 是否正确。

还可以把这些偏差除以本组奖励的标准差，统一比较尺度。
先算四个样本的总体标准差（分母为 $4$），再定义\textbf{组内相对优势} $\mathcal A_i$：
\[
\sigma_r=\sqrt{\tfrac14\sum_{i=1}^4(r_i-0.5)^2}=0.5,
\qquad \mathcal A_i=\frac{r_i-\bar r}{\sigma_r},
\]
因此组内相对优势为
\[
(\mathcal A_1,\mathcal A_2,\mathcal A_3,\mathcal A_4)=\boxed{(+1,-1,+1,-1)}.
\]
后续更新用正优势鼓励这次回答、用负优势抑制它；怎样沿梯度实现，见紧接着的选读算例。
减均值确定相对好坏，除标准差调整尺度；若整组奖励相同，$\sigma_r=0$，加入数值稳定项避免除零后，各个优势仍为零，这组样本的策略项不提供区分性梯度。
\end{example}

\subsubsection*{\texorpdfstring{{\symboliconfont ★}}{★} 选读：归一化口径与 logit 更新}
归一化口径因实现而异：改用样本标准差（分母为 $3$）时 $\sigma_r=\sqrt{1/3}\approx0.577$，
优势约为 $\pm0.866$；Dr.~GRPO 不做组内标准差归一化，因此本例优势为 $\pm0.5$，并移除了按回答长度归一化损失的因素。
比较报告时要写明所用口径。
\textbf{注意：}此处只算组内优势；完整 GRPO 还涉及策略概率比、裁剪和正则化。
可参见 \href{https://arxiv.org/html/2402.03300v3}{DeepSeekMath}
与 \href{https://arxiv.org/abs/2501.12948}{DeepSeek-R1}；
不除标准差的分析见 \href{https://arxiv.org/abs/2503.20783}{Liu 等（2025）}。

\noindent\textbf{优势怎样推高或压低一个 logit？}\par
把已采样答案与优势固定，先忽略裁剪和正则化，单个样本的策略损失写成
\[
\ell_i=-\mathcal A_i\ln p_\theta(y_i),\qquad
\frac{\partial\ell_i}{\partial Z_j}
=\mathcal A_i\bigl(p_j-\mathbf1_{j=y_i}\bigr).
\]
它就是第~\ref*{intro:20c-classification-loss}~章的 $p-\mathrm{onehot}$ 乘上优势：
正优势提高所选答案的分数，负优势压低它。
仅就单样本形式而言，$\mathcal A_i=+1$ 时得到 SFT 的交叉熵梯度；
SFT 的示范来自数据集，策略学习的动作来自采样、反馈来自奖励。

沿用 $(42,41,42,50)$ 与 $(+1,-1,+1,-1)$。
把整条回答暂作一个三分类动作，候选依次为 $42,41,50$，初始 $Z=(0,0,0)$、$p_j=1/3$。
四项梯度取平均为
\[
\begin{aligned}
\nabla_Z L
&=\tfrac14\bigl[(-2/3,1/3,1/3)+(-1/3,2/3,-1/3)\\
&\hspace{17mm}+(-2/3,1/3,1/3)+(-1/3,-1/3,2/3)\bigr]\\
&=\boxed{(-0.5,0.25,0.25)}.
\end{aligned}
\]
学习率 $\eta=0.2$ 时，一步得到 $Z'=(0.1,-0.05,-0.05)$，
正确答案的概率从 $1/3$ 升至 $\boxed{0.367456}$。
真实语言策略对回答中的各 Token 计分；完整 GRPO 还要组合旧策略概率比、裁剪和正则化。
此处隔离的是\textbf{奖励方向怎样沿 Softmax 梯度传回分数}。
脚本 \texttt{icl\_policy\_demo.py} 还用有限差分核对这三个导数。

把自动验证用在更复杂的任务上，还需要检查答案格式、代码测试结果、奖励漏洞以及验证器本身的错误。



\subsection{冻结参数，把预算分给多次尝试}
\phantomsection\label{concept:test-time-compute}

\textbf{把计算预算分给更多尝试，再用验证器筛选。}
前面的奖励进入优化器，改变未来请求使用的参数；现在冻结参数，
把同一个验证器用于本次候选的筛选。采样与验证只改变运行状态。

沿用 RLVR 的 $17+25$ 算例，验证器直接执行加法，正确答案是 $42$。
现在\textbf{冻结参数}，让模型生成 $n$ 份完整候选，再对每份运行验证器，
选出通过验证的答案；全部未通过就报告本次未解出。
这就是用完美二值验证器做 Best-of-$n$ 的最小版本。
一般的 Best-of-$n$ 也可按奖励模型的分数选最高项，其准确性取决于评分器。

\begin{example}[breakable=false]
\noindent\emph{从一次尝试到八次：先算全部失败的概率}\par
假设每次\textbf{独立}采样都从 $\{41,42,43,50\}$ 中等概率取一个答案，
答对的概率固定为 $p=0.25$。
一次失败概率是 $0.75$；$n$ 次全失败概率是 $0.75^n$。
因此至少有一份正确候选的概率为
\[
P_n=1-(1-p)^n.
\]
\[
\begin{array}{c|ccccc}
n&1&2&4&8&16\\\hline
P_n&0.2500&0.4375&0.6836&0.8999&0.9900
\end{array}
\]
完美验证器能选出存在的正确答案，所以这里的最终成功率也等于 $P_n$。
从 $4$ 次增到 $8$ 次提高约 $21.63$ 个百分点，
从 $8$ 次增到 $16$ 次只再提高约 $9.01$ 个百分点：预算继续翻倍，收益逐渐变小。
\end{example}

\par\medskip
\Needspace{85mm}
\begin{center}
\tagstructbegin{tag=Figure,alttext={Best-of-n 教学成功率图。横轴为一到十六份候选，纵轴为交付正确答案的概率。独立候选且完美核验时，单次成功率四分之一，曲线经过 n 为一、二、四、八、十六时的百分之二十五、四十三点七五、六十八点三六、八十九点九九、九十九点零零。候选完全复制时，成功率始终为百分之二十五。}}
\begin{tikzpicture}
\begin{axis}[width=.88\linewidth,height=61mm,
 xmin=1,xmax=16,ymin=0,ymax=1.06,
 xtick={1,2,4,8,16},ytick={0,.25,.5,.75,1},
 xlabel={候选数 $n$},ylabel={成功率},
 tick label style={font=\footnotesize},label style={font=\small},
 axis lines=left,axis line style={draw=muted},
 grid=major,grid style={draw=linegray},
 legend style={font=\footnotesize,draw=none,fill=white,at={(.97,.18)},anchor=south east}]
\addplot[datacolor,very thick,domain=1:16,samples=100]{1-.75^x};
\addlegendentry{独立采样＋完美核验}
\addplot[only marks,mark=*,datacolor] coordinates {(1,.25) (2,.4375) (4,.68359375) (8,.89988708496) (16,.98997740424)};
\addlegendentry{手算的五个点}
\addplot[noisecolor,thick,dashed,domain=1:16]{.25};
\addlegendentry{完全复制同一次答案}
\end{axis}
\end{tikzpicture}
\captionof{figure}{单次成功率 $p=0.25$、完美核验下，独立采样与复制同一答案的成功率对比}
\label{fig:best-of-n-success}
\tagstructend
\end{center}


\textbf{现在动手：至少需要多少次，才能达到 $90\%$？}
由 $1-0.75^n\geq0.9$，得到
\[
n\geq\frac{\ln0.1}{\ln0.75}\approx8.004,
\qquad n_{\min}=9.
\]
八次的 $89.9887\%$ 虽可四舍五入成 $90\%$，仍低于题设门槛；
九次为约 $92.49\%$。
若每份候选固定生成 $128$ 个 Token，完整采样八份就是 $1024$ 个输出 Token，
还要加上验证与输入处理开销；并行执行可缩短等待，总运算与资源占用照样计入。

\begin{boundarybox}
$1-(1-p)^n$ 描述\textbf{固定单次成功率、独立候选}中至少一次答对的概率。
若八份候选完全复制同一次随机答案，成功率仍只有 $25\%$。
若验证器误把与 $42$ 相差 $1$ 的 $41$、$43$ 也判为正确，并交付第一份通过的候选：
至少一份通过的概率是 $1-0.25^n$，而第一份通过者在 $41,42,43$ 中等可能，
\[
P^{\rm faulty}_n=\frac{1-0.25^n}{3}:\quad
P^{\rm faulty}_1=0.25,\ P^{\rm faulty}_2=0.3125,\ P^{\rm faulty}_4\approx0.3320,\ P^{\rm faulty}_8\approx0.3333.
\]
再多采样也停在约 $1/3$：“候选中存在正确解”与“最后交付正确解”要分开评测。
在开放任务中，模型自评分、程序测试与资料出处分别有自己的覆盖范围，
验证器需要独立检查。
\end{boundarybox}

RLVR 把验证奖励用于\textbf{更新参数}；这里的 Best-of-$n$ 用验证结果\textbf{筛选本次候选}。
串行修正和搜索也能使用预算，实际策略要在相同任务和预算下比较；
可参阅 \href{https://arxiv.org/abs/2408.03314}{Snell 等对测试时计算的研究}。
\textbf{复算入口：}\texttt{python3 scripts/test\_time\_compute\_demo.py}，
脚本枚举小样本空间，分别核对独立采样、完全相关采样与错误验证器的结果。


\subsection{“让模型写步骤”与“训练模型推理”不是一回事}
\phantomsection\label{concept:cot-prompt}

\textbf{写出中间步骤，为后续 Token 提供一个可读取的工作区。}
思维链提示（Chain-of-Thought, CoT）让模型先写计算过程，再给结果。
每个新 Token 都经过一次模型前向计算，后续生成又能读取此前写出的步骤，
于是多步任务得到更多串行计算与中间记录。

CoT 提示只改变输入，训练才修改参数，推理时搜索则额外花费计算；
三者的分工见本章开头的\hyperref[sec:training-inference-roles]{分工表}。

\begin{boundarybox}
步骤写得完整仍需逐项核对；它们是生成文本，无法直接证明模型内部采用了同一推理过程。
错误步骤会向后传播，额外生成也会增加 Token 与延迟成本。
训练后的收益用独立任务测试和失败案例衡量。
\end{boundarybox}

\paragraph{回到季度报告任务}
Few-shot 给格式示例，SFT 用核对过的示范训练格式与缺证据时的回应，偏好反馈比较来源与单位是否清楚，
可验证奖励检查数值、公式与出处。读取报告的权限与原表证据仍由应用提供；
\hyperref[intro:23d-agent-systems]{检索与任务系统}一章将接着说明证据怎样进入上下文。


\chapterchecktasks{23c-post-training}{理解、手算与判断}{
\label{check:23cc-test-time-compute}
\item\textbf{比较：}预训练、SFT 与 Few-shot ICL 中哪些更新参数？再比较 SFT、成对偏好学习与可验证奖励各收到什么反馈。\chapterinlinehint{先检查是否反传更新，再列示范、比较与核验结果。}
\item\textbf{手算：}SFT 目标 Token 的概率为 $0.25$，求单位置交叉熵；若四个组内奖励是 $(1,0,1,0)$，按总体标准差标准化得到哪些相对优势？\chapterinlinehint{先算负对数，再用奖励减均值并除标准差。}
\item\textbf{{\symboliconfont ★} 找错：}同学把这四个相对优势直接当成完整 GRPO 更新，遗漏了什么策略目标与数值边界？\chapterinlinehint{在相对优势之后继续写策略目标，并检查零方差。}
\item\textbf{解释：}向同学解释 RLVR 与 Best-of-$n$ 各把验证结果用在何处。\chapterinlinehint{检查本次是否更新模型权重。}
\item\textbf{估算：}若每份候选独立答对概率为 $1/2$，完美验证器检查四份，最终成功率是多少？\chapterinlinehint{先求所有候选都失败的概率。}
\item\textbf{复现：}写十行以内 Python，比较四份独立候选与四份完全相同候选的成功率。解释两者为什么不同。\chapterinlinehint{分别枚举独立样本空间与共享随机变量。}
}
\par\smallskip
训练改变参数；测试时计算把固定模型的预算分给生成、搜索和验证。
判断方法是否值得采用时，还要同时报告单次成功率、候选相关性、验证质量与总成本。\phantomsection\label{route:feedback-end}
